\begin{slide}
	%\vspace{-10mm}
	\centering
	\raisebox{-.5\height}{\includegraphics[scale=.8]{../problems/longexposure/solution/samplein}}
	\hspace{1mm}
	->
	\hspace{1mm}
	\raisebox{-.5\height}{
\includegraphics[scale=.8]{../problems/longexposure/solution/sampleans}}
	\begin{block}{Problem}
		A white surface with $n$ black rectangles is rotated around its origin fast.
		A photo is taken of the entire rotation.
		Determine the `black' surface area of the photo (parts covered by rectangles throughout the entire rotation)
		as well as the `grey' surface area (parts covered by rectangles through some but not all of the rotation).
	\end{block}
\end{slide}
\begin{slide}
	\centering
	\includegraphics[scale=1.25]{../problems/longexposure/solution/diags}
	\vspace{-47mm}
	\begin{block}{Solution (Variant 1)}
		\begin{itemize}
			\item As a circle's radius increases, the number of $1 \times 1$ squares intersected by its outline only increases for every new integer.
			\item Number of intersected $1 \times 1$ squares for radius $r$ is $8 \cdot \lfloor r \rfloor + 4$.
			\item Squares can be grouped by the `diagonal' they lie on.
		\end{itemize}
	\end{block}
\end{slide}
\begin{slide}
	\centering
	\includegraphics[scale=1.25]{../problems/longexposure/solution/intervals}
	\vspace{-47mm}
	\begin{block}{Solution (Variant 1)}
		\begin{itemize}
			\item Each rectangle is essentially a set of intervals, one interval per affected diagonal.
			\item For every integer radius, use the interval entry and exit points (up to the next integer)
			      to compare the covered diagonals to the required number of $1 \times 1$ squares.
			\item Add surface areas of the resulting black and grey `rings' to the respective total.
		\end{itemize}
	\end{block}
\end{slide}
\begin{slide}
	\begin{block}{Solution (Variant 2)}
		\begin{itemize}
			\item Do coordinate compression, creating an approximately $n \times n$ grid.
			\item Each rectangle covers some fields in the grid. Mark those as black. All unmarked fields are white. ($O(nwh)$)
			\item Fields reach minimal/maximal distance to the centre at a vertex or at x = 0 or y = 0. Thus, we can calculate the distance interval this field covers.
			\item Sort and merge black and white intervals ($O(n^2 \log(n))$ for $n^2$ intervals). If black and white intervals overlap, the overlapping part is grey.
		\end{itemize}
	\end{block}
	\pause
	\begin{block}{Optimization}
		The second step can be optimized to run in $O(n^2 \log(n))$ using a sweepline approach with a segment tree. This was not required.
	\end{block}
\end{slide} 
